\documentclass[../../../include/open-logic-section]{subfiles}

\begin{document}

\olfileid{sth}{card-arithmetic}{ch}

\olsection{连续统假设}

上一节的结果正确地提示了：如果无穷基数能保证与 $2$ 的幂“顺畅配合”，即（根据 \olref[opps]{lem:SizePowerset2Exp}）与取幂集操作顺畅配合，那么基数幂运算将会相当\emph{简易}。但我们不能假设无穷基数确实\emph{能}与幂集顺畅配合。

为了深入探讨这一点，我们引入一些简洁的记号。

\begin{defn}
令 $\cardsucc{\cardfont{a}}$ 为严格大于 $\cardfont{a}$ 的最小基数，我们定义如下两个无穷序列：
\begin{align*}
	\aleph_{0} &\defis \omega & 		
	\beth_{0} &\defis \omega\\
	\aleph_{\alpha \ordplus 1} &\defis \cardsucc{(\aleph_{\alpha})} &
	\beth_{\alpha+1} &\defis \cardexpo{2}{\beth_{\alpha}}\\
	\aleph_{\alpha} &\defis \bigcup_{\beta< \alpha} \aleph_{\beta} &
	\beth_{\alpha} &\defis \bigcup_{\beta < \alpha}\beth_{\beta} & \text{当 $\alpha$ 是极限序数时}.
\end{align*}
\end{defn}

$\cardsucc{\cardfont{a}}$ 的定义是妥当的，因为 \olref[cardinals][classing]{lem:NoLargestCardinal} 告诉我们，对每个基数 $\cardfont{a}$，都存在某个大于 $\cardfont{a}$ 的基数，而超限归纳法保证了存在一个大于 $\cardfont{a}$ 的\emph{最小}基数。该定义的其余部分则由超限递归给出。

Cantor（康托尔）引入了这个“$\aleph$”记号；这是 \emph{aleph}，希伯来字母表的第一个字母，也是希伯来语中“无穷”一词的首字母。Peirce（皮尔士）引入了“$\beth$”记号；这是 \emph{beth}，希伯来字母表的第二个字母。\footnote{Peirce 在 1900 年 12 月写给 Cantor 的信中使用了这个记号。不幸的是，Peirce 在那里还给出了一个错误的论证，声称对于 $\alpha \geq \omega$，$\beth_\alpha$ 不存在。}由此，这些记号便为我们给出了无穷基数。

\begin{prop}
对所有序数 $\alpha$，$\aleph_\alpha$ 和 $\beth_\alpha$ 均为基数。
\end{prop}

\begin{proof}
两个结论通过简单的超限归纳即可证明。$\aleph_0 = \beth_0 = \omega$ 是基数，由 \olref[cardinals][classing]{omegaisacardinal} 可知。假设 $\aleph_\alpha$ 和 $\beth_\alpha$ 都是基数，则 $\aleph_{\alpha+1}$ 和 $\beth_{\alpha+1}$ 由显式定义即为基数。而根据 \olref[cardinals][classing]{unioncardinalscardinal}，一组基数的并集仍是基数。
\end{proof}
\noindent
此外，每一个无穷基数都是一个 $\aleph$：

\begin{prop}
若 $\cardfont{a}$ 是无穷基数，则存在唯一的 $\gamma$ 使得 $\cardfont{a} = \aleph_\gamma$。
\end{prop}

\begin{proof}
对基数进行超限归纳。归纳假设：若 $\cardfont{b} < \cardfont{a}$ 则 $\cardfont{b} =
\aleph_{\gamma_\cardfont{b}}$。若存在某个 $\cardfont{b}$ 使得 $\cardfont{a} =
\cardsucc{\cardfont{b}}$，则 $\cardfont{a} = \cardsucc{(\aleph_{\gamma_\cardfont{b}})}= \aleph_{\gamma_\cardfont{b}+1}$。若 $\cardfont{a}$ 不是任何基数的后继，那么由于基数是序数，$\cardfont{a} = \bigcup_{\cardfont{b} < \cardfont{a}} \cardfont{b} =
\bigcup_{\cardfont{b} < \cardfont{a}}{\aleph_{\gamma_\cardfont{b}}}$，因此 $\cardfont{a} = \aleph_\gamma$，其中 $\gamma =
\bigcup_{\cardfont{b} < \cardfont{a}}\gamma_\cardfont{b}$。 
\end{proof}

既然每个无穷基数都是一个 $\aleph$，这促使我们发问：是否每个无穷基数也都是一个~$\beth$？当然，\emph{倘若}果真如此，那么无穷基数将与取幂集操作“顺畅配合”。事实上，我们将得到以下结论：

\begin{defish}
\emph{广义连续统假设} (GCH)。对所有 $\alpha$，$\aleph_\alpha  = \beth_\alpha$。
\end{defish}

此外，倘若 GCH 成立，我们就能对基数幂运算做相当程度的简化。特别地，当 $\cardfont{b} < \cardfont{a}$ 时，我们可以证明 $\cardexpo{\cardfont{a}}{\cardfont{b}}$ 的值会被 $\cardfont{a}\leq \cardexpo{\cardfont{a}}{\cardfont{b}} \leq
\cardsucc{\cardfont{a}}$ 所限定。进而，我们可以给出精确的条件，以判定这两种可能性中哪一个成立（即究竟是 $\cardfont{a} = \cardexpo{\cardfont{a}}{\cardfont{b}}$ 还是 $\cardexpo{\cardfont{a}}{\cardfont{b}} =
\cardsucc{\cardfont{a}}$）。\footnote{该条件由\emph{共尾度}决定。}

但 GCH 是一个\emph{假说}，而非一个\emph{定理}。事实上，\citet{Godel1938} 证明了如果 $\ZFC$ 是一致的，那么 $\ZFC + \text{GCH}$ 也是一致的。但后来发现，我们同样可以将 $\lnot$GCH 加到 $\ZFC$ 中。考虑 GCH 最简单且非平凡的\emph{实例}，即：

\begin{defish}
\emph{连续统假设} (CH)。$\aleph_1 = \beth_1$。
\end{defish}

\citet{Cohen1963} 证明了如果 $\ZFC$ 是一致的，那么 $\ZFC + \lnot\text{CH}$ 也是一致的。因此连续统假设独立于 $\ZFC$。

该假设之所以被称为“连续统假设”，是因为“连续统”是实数轴 $\Real$ 的另一个名称。\olref[opps]{continuumis2aleph0} 告诉我们 $\card{\Real} = \beth_1$。所以连续统假设说的是：不存在其他基数介于自然数的势 $\aleph_0 = \beth_0$ 与连续统的势 $\beth_1$ 之间。

既然（广义）连续统假设独立于 $\ZFC$，我们该如何看待它们的\emph{真值}呢？能说的\emph{太多}了。事实上，集合论中许多富有成果的近期研究都致力于探讨这些问题。但有两点非常简短的看法，无疑值得强调。

第一：从这些形式上的独立性结果，并不能\emph{直接}推出广义连续统假设或连续统假设的真值是\emph{不确定的}。毕竟，也许我们只需要添加更多在我们看来很自然的公理，就能非此即彼地解决这一问题。G\"odel（哥德尔）本人就曾暗示这是正确的应对方式。

第二：CH 独立于 $\ZFC$ 这个结果当然是\emph{惊人的}，但它绝不是\emph{不可信的}（就其字面意义而言）。这一点很简单：就 $\ZFC$ 所揭示的信息而言，从一个基数过渡到它的后继，所用的工具可能比简单地取幂集更精细。
 %The operation of taking powersets moves us from one stage of the
 %hierarchy to its successor stage (from $V_{\alpha}$ to
 %$V_{\alpha+1}$). 

给出了这两点观察之后，如果你想了解更多，现在就得去向那些在此问题上各执己见的哲学家和数学家求教了。\footnote{不过，你可能想从阅读 \citet[\S15.6]{Potter2004} 开始。}

\end{document}